\chapter{살아 있는 뉴런으로 만든 기계}
\chaptermark{살아 있는 뉴런으로 만든 기계}
\label{sec:livingmachine}

\section{회로도가 보이는 테스트 벤치}

\ref{sec:debugging}장에서 우리는 회로도가 없는 기계를 디버깅했다. 프로브는 800억 개 뉴런 가운데 천 개만 들을 수 있고, 나머지는 모두 추측해야 한다. 이런 처지의 엔지니어라면 간단한 일을 할 것이다. 회로의 작은 조각을 브레드보드 위에 조립해서 부품 하나하나가 보이게 하는 것이다. 뉴런으로도 이것을 할 수 있다.

피질 뉴런을 떼어 내어 전극 배열이 있는 칩 위에서 키운다. 전극은 수십 개에서 수만 개까지다. 몇 주가 지나면 세포는 돌기를 뻗어 수천에서 수십만 개 뉴런의 네트워크로 연결된다. 대부분은 \nt{GLUT}이고 \nt{GABA}는 더 적은데, 그 비율은 표본에 따라 다르다. 해마 배양에서는 뉴런의 약 $6\%$이다~\cite{Benson1994}. 각 전극은 \ref{sec:debugging}장의 프로브처럼 들을 수도 있고, 테스트 신호처럼 전류를 넣을 수도 있다. 두 번째 종류의 테스트 벤치는 \emph{오가노이드}다. 인간 줄기세포에서 키운 지름 1밀리미터 정도의 3차원 신경 조직 덩어리이다~\cite{Lancaster2013}. 이런 테스트 벤치에서 살아 있는 네트워크는 몇 달 동안 작동한다. 오가노이드 실험을 네트워크로 원격 수행하면서 100일 넘게 연속으로 돌릴 수 있는 시설도 있다~\cite{Jordan2024}. 이 분야를 \emph{생물학적 컴퓨팅} 또는 \emph{오가노이드 지능}이라고 부른다~\cite{Smirnova2023}.

테스트 벤치는 뇌와 두 가지 중요한 점에서 다르다. 우선 아주 작다. 뉴런 수가 10만 배 적다. 그리고 브로드캐스트 채널이 없다. 피질 배양에는 \nt{DOPA}, \nt{SERO}, \nt{NORA}, \nt{ACET} 뉴런이 없다. 데이터 버스와 흐름 제어는 있지만 제어 레지스터는 없는 기계이다. 이 기계가 무엇을 할 수 있는지 보자.

\section{네트워크가 스스로 하는 일}

가만히 두어도 배양은 조용하지 않고, 고르게 잡음을 내지도 않는다. 이따금 네트워크 거의 전체가 한꺼번에 발화하는 \emph{일제 발화}가 일어나고, 다시 잠잠해진다. 일제 발화 사이의 간격은 배양에 따라 1초에서 몇 분이다~\cite{Wagenaar2006}. 엔지니어에게는 익숙한 그림이다. 강한 양의 되먹임을 걸고 부하를 떼어 낸 증폭기는 발진기가 된다.

메커니즘은 이미 알고 있다. \nt{GLUT} 뉴런 사이의 강한 상호 연결은 사태를 일으킨다. \ref{sec:gaba}장에서 명제~\ref{prop:ei}의 조건이 깨졌을 때와 같다. 사태는 시냅스의 신경전달물질 저장량을 빠르게 고갈시키고~\eqref{eq:depression}, 네트워크는 조용해진다. 저장량은 시간 상수 $\tau_{\text{rec}}$로 회복되고, 새 사태에 충분한 비율 $x^*$까지 회복되면 다음 일제 발화가 온다. 일제 발화 간격은
\begin{empheq}[box=\keybox]{equation}
T_{\text{burst}} \;\approx\; \tau_{\text{rec}}\,\ln\frac{1}{1-x^*} .
\label{eq:burstinterval}
\end{empheq}
\ref{sec:code}장의 $\tau_{\text{rec}}=0.8\,\text{s}$와 $x^*=0.9$를 넣으면 약 2초, $x^*=0.99$이면 약 4초가 나온다. 이것은 이 책의 $\tau_{\text{rec}}$를 쓴 모델의 추정값이며, 측정된 간격의 짧은 쪽 끝에 들어맞는다. 미세 배양에서 직접 측정해 보면 일제 발화 뒤 신경전달물질 저장량이 회복되는 데 수십 초가 걸린다~\cite{CohenSegal2011}. 즉 거기서는 $\tau_{\text{rec}}$가 더 크다. 그런 $\tau_{\text{rec}}$를 쓰면 식~\eqref{eq:burstinterval}은 느린 배양에서처럼 수십 초에서 몇 분을 준다. 이것은 이완형 발진기이며, \ref{sec:muscles}장의 리듬 발생기와 친척이다. 거기서는 뉴런 무리의 피로가 발진기를 다시 충전했고, 여기서는 신경전달물질 저장량이 그 일을 한다.

일제 발화는 네트워크가 할 일이 없을 때 하는 일이다. 살아 있는 뇌에서도 비슷한 그림이 깊은 잠(\ref{sec:sleep}장)과, 진짜 입력이 들어오기 전의 발달 중인 뇌에서 나타난다. 이 유사성은 시사하는 바가 있지만, 메커니즘이 같다는 것은 가설이다.

\section{도파민 교사 없이 네트워크를 가르치는 법}

배양에는 \nt{DOPA}가 없으므로 “더 좋다, 더 나쁘다”라는 신호는 우리가 직접 전극으로 넣어야 한다. 실험을 보면 테스트 벤치 위의 네트워크는 실제로 응답을 바꾼다. 그리고 가장 흥미로운 것은 \emph{무엇으로} 가르치느냐이다.

\emph{입을 다무는 교사.} 한 전극을 통해 1--3초마다 한 번씩 네트워크를 자극하다가, 다른 전극에서 자극 후 $50\pm10\ms$에 원하는 응답이 나타나면 자극을 5분 동안 끈다. 이런 주기를 몇 번 거치면 원하는 응답이 자극에 곧바로 나타난다~\cite{Shahaf2001}. 여기서 보상은 침묵이다. 자극은 네트워크를 흔들고, 자극을 멈추면 네트워크는 올바른 응답을 낸 상태에 머문다. 이것은 \ref{sec:dofa}장의 무작위 섭동을 통한 경사 하강(명제~\ref{prop:gradient})이며, 여기서는 흔드는 것 자체가 오차 역할을 한다. 맞을 때까지 흔든다.

\emph{테니스를 치는 네트워크.} 전극이 촘촘히 배열된 테스트 벤치 실험에서, 약 100만 개의 뉴런을 라켓이 하나인 컴퓨터 테니스 게임에 연결했다~\cite{Kagan2022}. 공의 위치는 “감각” 전극에 전류로 넣었다. 공이 어디 있는지는 장소로, 얼마나 먼지는 빈도로 부호화했다. 두 “운동” 영역의 활동이 라켓을 위아래로 움직였다. 학습 규칙은 특이했다. 라켓이 공을 받아 치면 네트워크는 짧고 \emph{예측 가능한} 자극을 받았고, 놓치면 몇 초 동안 \emph{예측 불가능한} 무작위 전류를 받았다. 20분 동안 게임하는 사이에 랠리가 길어졌다. 세션 안에서 유의미한 향상은 완전한 되먹임일 때만 있었다. 자극 대신 침묵을 준 경우에도 배양은 세션 끝에 되먹임이 없을 때보다 더 잘 쳤지만 효과가 작았고, 며칠에 걸친 세션 간의 안정적인 학습은 없었다. 이 실험의 자세한 분석은 \ref{sec:dishbrain}장에 있다.

저자들은 네트워크가 자기 입력이 예측 가능해지도록 행동한다는 가설로 결과를 설명했다~\cite{Friston2010}. \ref{sec:code}장의 스파이크 절약, \ref{sec:recognition}장의 프레임 예측과 같은 생각이다. 입력이 예상대로일 때 기계는 덜 소비한다. 그러나 실험 자체도, 특히 배양의 “지각력”에 관한 저자들의 표현도 날카로운 비판을 받았다. 효과가 작고 더 단순하게 설명할 수도 있다는 것이다~\cite{Balci2023}. 정직하게 평가하면 이렇다. 테스트 벤치의 배양이 되먹임에 따라 행동을 바꾼다는 것은 측정된 사실이고, 그것이 바로 자기 입력을 예측하도록 학습한다는 것은 가설이다.

\emph{몸을 움직이는 네트워크.} 비슷한 방식으로 배양을 단순한 가상 “몸”에 연결하고, 훈련 자극의 시공간 패턴을 골라 몸이 목표로 움직이도록 가르쳤다~\cite{Bakkum2008}. 이 실험들 어디에도 도파민은 없다. 교사는 엔지니어가 고른 전류 패턴이다. 교사가 프로그램이나 도파민 자체가 되면 무엇이 달라지는지는 \ref{sec:dishdopamine}--\ref{sec:cartpole}절에서 다룬다.

\begin{figure}[t]
\centering
\begin{tikzpicture}[>=Stealth,
  blk/.style={draw,very thick,rounded corners=2pt,align=center,font=\footnotesize,minimum height=0.8cm,fill=blue!5}]
% (a) closed loop
\node[font=\small] at (1.5,4.3) {(a) 닫힌 루프};
\node[blk,text width=2.6cm] (G) at (1.5,3.3) {게임: 공과 라켓};
\draw[very thick,fill=orange!10,rounded corners=3pt] (0.5,0.2) rectangle (2.5,1.6);
\foreach \x in {0.7,1.0,...,2.35}{\foreach \y in {0.4,0.7,1.0,1.3}{\fill[gray!60] (\x,\y) circle (0.04);}}
\draw[->,thick,blue!60!black] (G.west) -| (-0.5,0.9) -- (0.5,0.9);
\node[font=\scriptsize,blue!60!black,anchor=east] at (-0.6,2.1) {공 $\to$ 전류};
\draw[->,thick,red!70!black] (2.5,0.9) -- (3.5,0.9) |- (G.east);
\node[font=\scriptsize,red!70!black,anchor=west] at (3.6,2.1) {스파이크 $\to$ 라켓};
\node[font=\scriptsize] at (1.5,-0.2) {전극 배열 위의 뉴런};
\node[font=\scriptsize,align=center] at (1.5,-0.85) {히트: 예측 가능한 자극\\미스: 무작위 전류};
% (b) reservoir
\begin{scope}[xshift=7.9cm]
\node[font=\small] at (1.6,4.3) {(b) 저수지};
\node[blk,text width=1.1cm] (X) at (-0.4,1.0) {입력\\$u(t)$};
\draw[very thick,fill=orange!10] (1.6,1.0) circle (0.8);
\foreach \a/\r in {20/0.35,80/0.5,150/0.3,210/0.55,280/0.4,330/0.2,110/0.15}{\fill[red!60!black] ({1.6+\r*cos(\a)},{1.0+\r*sin(\a)}) circle (0.05);}
\node[font=\scriptsize] at (1.6,-0.2) {오가노이드: 교사 없음};
\node[blk,text width=1.8cm,fill=green!8] (W) at (4.0,1.0) {판독\\$W_{\text{out}}$};
\draw[->,thick] (X) -- (0.8,1.0);
\draw[->,thick] (2.4,1.0) -- node[above,font=\scriptsize] {$\mathbf r(t)$} (W);
\draw[->,thick] (W) -- (4.0,2.3) node[above,font=\scriptsize] {$y(t)$};
\node[font=\scriptsize,green!40!black] at (4.0,0.3) {교사 있는 학습};
\end{scope}
\end{tikzpicture}
\caption{살아 있는 네트워크가 계산하게 만드는 두 가지 방법. (a) 닫힌 루프: 공의 위치를 전류로 넣고, 운동 영역의 스파이크가 라켓을 움직이며, 타격 뒤의 예측 가능한 자극 또는 무작위 자극이 교사 역할을 한다. (b) 저수지~\eqref{eq:reservoir}: 오가노이드에는 학습 신호를 주지 않는다. 오가노이드는 입력을 기억을 가진 수많은 비선형 응답으로 쪼개고, 오차로 학습하는 것은 바깥의 선형 판독 하나뿐이다. 이때 오가노이드 자체의 연결도 교사 없이 바뀐다.}
\label{fig:livingmachine}
\end{figure}

\section{살아 있는 저수지}

살아 있는 네트워크를 쓰는 다른 방법도 있다. 아예 가르치지 않는 것이다. 공학에서는 이것을 \emph{저수지 컴퓨팅}(reservoir computing)이라고 부른다~\cite{Maass2002,JaegerHaas2004}. 기억을 가진 기성 비선형 시스템을 가져온다. 입력에 대한 응답이 풍부하고 최근 과거에 의존하기만 하면 무엇이든 된다. 입력 $u(t)$가 시스템을 흔들면 응답 벡터 $\mathbf r(t)$가 나오고, 학습하는 것은 선형 판독뿐이다.
\begin{empheq}[box=\keybox]{equation}
y(t) \;=\; W_{\text{out}}\,\mathbf r(t) ,
\label{eq:reservoir}
\end{empheq}
그 가중치는 \ref{sec:motorlearning}장의 최소제곱 규칙~\eqref{eq:lms}으로 맞춘다. 저수지는 \ref{sec:motorlearning}장의 작은 \nt{GLUT} 뉴런들이 하던 일을 한다. 입력을 긴 벡터로 쪼개어, 필요한 클래스들이 평면으로 분리될 수 있게 만든다(\ref{sec:dendrites}장).

전극 배열 위의 오가노이드는 이런 저수지 역할에 알맞았다. 전류 패턴으로 넣은 말소리를 선형 판독 하나만 학습시킨 뒤 화자별로 구분했더니 정확도가 약 $78\%$였다~\cite{Cai2023}. $78\%$라는 정확도는 저자들이 직접 보고하고 언론이 되풀이한 숫자이다. 쓸모 있는 결과지만, 그 안에서 “지능”이 어디에 있는지 이해하는 것이 중요하다. 오가노이드는 비선형성과 감쇠하는 기억을 제공하고, 오차에 의한 학습은 바깥의 실리콘에서 이루어진다(그림~\ref{fig:livingmachine}). 그렇다고 오가노이드가 그대로인 것은 아니다. 저자들에 따르면 훈련 동안 오가노이드 자체의 기능적 연결성이 바뀌었고, 이를 오가노이드 안에서 일어나는 교사 없는 학습이라고 부른다~\cite{Cai2023}. 정확도 가운데 얼마가 이 변화에서 오고 얼마가 판독에서 오는지는 분리되지 않았다.

\section{테스트 벤치가 기계에 대해 말해 주는 것}

\emph{에너지.} 추정~\eqref{eq:energy}에 따르면 스파이크 하나는 약 $10^{-10}$~J이 든다. 뉴런 100만 개의 배양이 초당 스파이크 하나의 빈도로 발화하면 신호 전달에 쓰는 전력은
\begin{empheq}[box=\keybox]{equation}
P \;\approx\; 10^{6}\times1\,\text{s}^{-1}\times10^{-10}\,\text{J} \;=\; 10^{-4}\,\text{W},
\label{eq:dishpower}
\end{empheq}
즉 0.1밀리와트이다. 이 스파이크를 자극하고 기록하고 처리하는 전자 장치는 몇 자릿수 더 많이 쓴다. \ref{sec:debugging}장에서처럼 병목은 계산이 아니라 계산과의 통신이다.

\emph{빠진 부품.} 테스트 벤치는 제어 레지스터 없이 \nt{GLUT} 버스와 \nt{GABA} 흐름 제어가 얼마나 많은 일을 할 수 있는지 보여 준다. 스스로 조직되고, 일제 발화를 만들고, 되먹임에 따라 응답을 바꾸고, 좋은 저수지 역할을 한다. 제어 레지스터 없이 할 수 없는 것은 무엇을 성공으로 볼지 스스로 정하는 일이다. 테스트 벤치에서는 엔지니어가 이 신호를 넣고, 뇌에서는 \ref{sec:dofa}--\ref{sec:acet}장에서 본 것처럼 브로드캐스트 채널이 넣는다.

\emph{윤리.} 오가노이드는 인간 세포에서 키우며, 복잡해질수록 그런 조직이 무언가를 느낄 수 있느냐는 질문이 무거워진다. 공학적 관점은 부분적인 답을 준다. \ref{sec:pain}장의 통증은 센서 하나의 신호가 아니라 관문, 음량 조절기, 브로드캐스트 채널을 갖춘 경보 시스템 전체이며, 테스트 벤치에는 그런 것이 없다. 그러나 이것은 추론이지 증명이 아니다. 이 분야 스스로도 처음부터 실험과 함께 윤리적 평가를 진행하자고 제안한다~\cite{Smirnova2023}.

\begin{keyblock}
\begin{proposition}[테스트 벤치 위의 살아 있는 네트워크]
\label{prop:livingmachine}
전극 배열 위의 \nt{GLUT} 뉴런과 \nt{GABA} 뉴런 네트워크는 스스로 일제 발화를 만들고~\eqref{eq:burstinterval}, 되먹임에 따라 응답을 바꾸며, 바깥에서 학습된 판독을 위한 저수지~\eqref{eq:reservoir} 역할을 할 수 있다. 이것은 모두 측정되었다. 이런 네트워크가 어떤 일반적 규칙, 예를 들어 자기 입력을 예측하는 규칙에 따라 학습한다는 것은 가설이며, 그 “지각력”에 관한 과장된 표현은 대다수 전문가가 받아들이지 않는다.
\end{proposition}
\end{keyblock}

\section{빛의 섬광으로 주는 도파민}
\label{sec:dishdopamine}

배양에 도파민을 교사로 준 직접 실험으로 우리가 아는 것은 하나뿐이다. 연구자들이 원격으로 접속하는 오가노이드 플랫폼에서 수행되었다~\cite{Jordan2024}. 오가노이드를 적시는 액체에 빛에 반응하는 “케이지”에 갇힌 도파민을 넣었다. 묶여 있는 분자는 수용체에 작용하지 않는다. 케이지 도파민의 농도는 $1$~mM이다. 네트워크에 보상을 주어야 할 때 자외선을 비추면 케이지가 분해되고 도파민이 부피 속으로 풀려난다.

루프는 이렇게 짜여 있다. 같은 자극을 되풀이해서 넣고, 그 자극이 이전보다 많은 스파이크를 일으키면 섬광을 켜서 도파민을 방출한다. $240$번 반복하는 동안 유발 스파이크 수는 전체적으로 늘었다. 저자들 스스로 이런 실험 하나만으로는 증가가 도파민 때문이라고 결론 내릴 수 없다고 쓴다~\cite{Jordan2024}.

엔지니어에게 이것은 3요소 규칙~\eqref{eq:threefactor}을 접시 안에 조립하려는 첫 시도이다. 송신자와 수신자의 활동이 흔적을 남기고, 공통 신호가 변화를 굳힐지 결정한다. 그러나 여기서 신호는 양수뿐이다. 보상이 있거나 없거나이며, 장치는 음의 오차 $\delta<0$를 내지 못한다. 게다가 도파민은 \eqref{eq:lowpass}의 농도처럼 퍼지고 천천히 제거되므로, 섬광이 잦으면 그저 전체 수준만 올라갈 수 있다. 학습을 이것과 구별하려면 대조 실험 두 가지가 필요하다. 네트워크의 응답과 무관한 무작위 시점에 같은 수의 섬광을 주는 것, 그리고 케이지 도파민 없이 같은 섬광을 주는 것이다. 또 휴식 뒤의 확인도 필요하다. 도파민이 사라진 뒤에도 변화가 남았는가?

\section{도파민이 실리콘을 가르친다}
\label{sec:biohybrid}

반대 방향의 실험에서는 살아 있는 세포가 전자 소자를 가르친다~\cite{Keene2020}. 도파민을 분비하는 세포를 유기 뉴로모픽 소자 위의 미세 채널에서 키운다. 이 소자는 트랜지스터이며, 채널의 전도도가 시냅스 가중치 역할을 한다. 세포에서 나온 도파민이 소자의 전극에서 산화되면서 전도도를 바꾼다. 이 장치는 틈새에서 신경전달물질을 치우는 메커니즘도 재현하므로, 가중치를 오래 바꿀 수 있을 뿐 아니라 원래대로 되돌릴 수도 있다.

회로로 보면 이것은 화학 입력을 받는 누설 적분기이다. 가중치는 도파민 흐름을 누적하고, “제거”가 얼마나 빨리 잊는지를 정한다. \eqref{eq:lowpass}와 같은 RC 회로이다. 중요한 것은 연결 방향이다. \ref{sec:debugging}장의 프로브는 브로드캐스트 레지스터를 보지 못한다. 그것이 화학적이기 때문이다. 이 소자는 바로 화학 레지스터를 읽어서 전기적 가중치로 바꾼다. 뇌-컴퓨터 인터페이스에게 이것은 데이터 버스뿐 아니라 제어 레지스터까지 듣는 센서로 가는 길이다.

\section{자기 도파민 뉴런을 가진 오가노이드}
\label{sec:midbrainorg}

인간 줄기세포로 \nt{DOPA} 뉴런이 있는 뇌 영역을 닮은 오가노이드를 키울 수 있다. 그 안에는 도파민을 분비하는, 작동하는 도파민 뉴런이 있다~\cite{Jo2016}. 이런 오가노이드는 주로 질병 연구를 위해 키운다. 그 도파민이 다른 네트워크의 교사 역할을 한 실험은 찾지 못했다.

살아 있는 뉴런으로 만든 기계에게 이것은 빠진 부품이다. \ref{sec:livingmachine}장의 테스트 벤치는 제어 레지스터 없는 데이터 버스와 흐름 제어이다. 여기에는 브로드캐스트 신호의 원천이 이미 있다. 모자란 것은 비평가, 즉 예측 오차~\eqref{eq:ctd}를 계산하고 이 뉴런들을 조종하는 네트워크이다. 피질 오가노이드를 이런 오가노이드와 연결해서 오차에 따라 도파민이 분비되게 하는 것은 열린 문제이며, 아직은 실험이 아니라 가설이다.

\section{도파민 없이 막대의 균형을 잡는 오가노이드}
\label{sec:cartpole}

최근 실험 가운데 가장 완전한 것은 도파민을 전혀 쓰지 않는다~\cite{Robbins2026}. 생쥐 피질 오가노이드를 “카트 위의 막대” 과제에 닫힌 루프로 연결했다. 오가노이드의 활동이 카트를 움직이고, 막대의 위치는 자극으로 되돌아 들어간다. 시도와 시도 사이에 오가노이드는 짧은 고주파 학습 자극을 받는다. 어떤 자극을 줄지는 강화 학습 프로그램이 고른다.

측정된 결과는 다음과 같다.
\begin{itemize}
\item 대부분의 오가노이드에서 프로그램이 고른 학습 신호가 무작위로 고른 신호나 신호가 없을 때보다 좋은 결과를 냈다.
\item $45$분 휴식 뒤에는 향상이 남지 않았다.
\item 두 종류의 글루탐산 수용체 AMPA와 NMDA를 모두 차단하면 향상이 사라졌다.
\end{itemize}

마지막 항목은 학습된 것이 어디에 사는지 말해 준다. \nt{GLUT} 시냅스 안이며, \ref{sec:weightstore}절에서 입력과 출력의 동시성 검출기로 작동하는 수용체를 통해서이다. 두 번째 항목은 무엇이 모자란지 보여 준다. 빠른 변화는 있지만 굳히기는 없다. \ref{sec:motorlearning}장에서 느린 학습자 없이 빠른 학습자만 있을 때와 같다. 교사의 역할도 바뀌었다. 전류 패턴을 고르는 엔지니어 대신 프로그램이 들어섰지만, 교사는 여전히 접시 바깥에 있다.

배양을 전극 배열에서 학습시킨 더 이른 실험들 중에서도 도파민을 쓴 것은 하나도 찾지 못했다. 학습 신호는 모두 전기 자극이었다.

\section{누가 배양을 가르치는가}

\begin{table}[t]
\centering
\footnotesize
\setlength{\tabcolsep}{4pt}
\renewcommand{\arraystretch}{1.15}
\begin{tabular}{>{\raggedright}p{2.7cm}>{\raggedright}p{2.5cm}>{\raggedright}p{3.2cm}>{\raggedright\arraybackslash}p{4.4cm}}
\toprule
실험 & 가르치는 주체 & 학습 신호 & 알려진 것 \\
\midrule
원하는 응답에서 자극 멈추기 & 엔지니어 & 자극의 종료 & 응답이 굳어진다 --- 측정됨 \\
DishBrain (\ref{sec:dishbrain}장) & 엔지니어 & 예측 가능한 전류 또는 무작위 전류 & 세션 안의 유의미한 랠리 증가는 완전한 되먹임일 때만, 침묵일 때는 효과가 더 작다 --- 측정됨. 세션 간에는 불안정. 원인은 논쟁 중 \\
카트 위의 막대 & 강화 학습 프로그램 & 프로그램이 고른 고주파 자극 & 무작위보다 낫지만 휴식 뒤에는 남지 않는다 --- 측정됨 \\
섬광으로 주는 도파민 & “스파이크가 늘면 보상” 규칙 & 빛으로 풀려난 도파민 & 스파이크 증가 --- 관찰. 도파민의 역할은 증명되지 않음 \\
바이오하이브리드 시냅스 & 살아 있는 세포 & 세포에서 나온 도파민 & 소자 가중치의 장기 변화와 복귀 --- 측정됨 \\
\nt{DOPA} 뉴런을 가진 오가노이드 & --- & --- & 도파민 원천은 있다. 교사로 쓰인 적은 없다 \\
\bottomrule
\end{tabular}
\caption{누가 무엇으로 칩 위의 살아 있는 뉴런을 가르치는가.}
\label{tab:dishteachers}
\end{table}

배양 자체가 학습하는 모든 실험에서 교사는 아직 바깥에 있다(표~\ref{tab:dishteachers}). 엔지니어, 프로그램, 또는 엔지니어가 정한 규칙이다. 도파민이 접시 안에 들어온 것은 단 한 번뿐이고, 그 역할은 아직 증명해야 한다. \ref{sec:dofa}장에서처럼 비평가, 오차, 흔적을 갖춘 진짜 브로드캐스트 채널을 접시 안에 조립하는 것은 아직 아무도 내딛지 않은 다음 걸음이다.


\section{정리}

\begin{table}[t]
\centering
\footnotesize
\setlength{\tabcolsep}{4pt}
\renewcommand{\arraystretch}{1.15}
\begin{tabular}{>{\raggedright}p{2.8cm}>{\raggedright}p{4.2cm}>{\raggedright\arraybackslash}p{5.8cm}}
\toprule
실험 & 살아 있는 네트워크가 하는 일 & 알려진 것 \\
\midrule
입력 없는 네트워크 & 1초에서 몇 분 간격의 일제 발화~\eqref{eq:burstinterval} & 측정됨. 시냅스 고갈을 통한 메커니즘은 널리 받아들여진 모델 \\
입을 다무는 교사 & 자극을 끄면 원하는 응답이 굳어진다 & 측정됨 \\
테니스 게임 & 랠리가 길어진다. 세션 안의 유의미한 향상은 완전한 되먹임일 때만 & 행동 변화는 측정됨. 세션 간 학습은 불안정. 효과의 크기와 입력 예측이라는 설명은 논쟁 중 \\
가상의 몸 & 알맞게 고른 자극 패턴으로 목표를 향해 움직인다 & 측정됨 \\
저수지 & 판독~\eqref{eq:reservoir}을 위한 비선형성과 기억 & 측정됨. 오차에 의한 학습은 바깥의 실리콘에서. 오가노이드의 연결도 바뀐다 \\
에너지 & 뉴런 100만 개당 약 $10^{-4}$~W~\eqref{eq:dishpower} & \eqref{eq:energy}에 따른 추정 \\
\bottomrule
\end{tabular}
\caption{살아 있는 뉴런으로 만든 기계가 보여 준 것.}
\label{tab:livingmachine}
\end{table}

살아 있는 뉴런으로 만든 기계(표~\ref{tab:livingmachine})는 제어 레지스터 없이 데이터 버스와 흐름 제어가 무엇을 할 수 있는지 보여 주는 테스트 벤치이다. 이들은 많은 것을 할 수 있지만, 무엇이 좋은지를 스스로 정하지는 못한다. 테스트 벤치에서는 자기 전류 패턴을 가진 엔지니어가 이 역할을 하고, 뇌에서는 이 책의 중간 부분에서 다룬 몇 안 되는 브로드캐스트 전선이 한다.

\section{연습문제}

\begin{exercise}[\lvlA]
\label{ex:21:1}
\emph{일제 발화 간격.} 일제 발화가 저장량을 0까지 고갈시키고, 그 뒤 $\tau_{\text{rec}}\,\dd x/\dd t=1-x$이며, $x=x^*$에서 새 일제 발화가 시작된다. $\tau_{\text{rec}}=0.8$~s. (a)~$x^*=0.9$와 $0.99$일 때의 간격을 구하라. (b)~문턱값 $x^*$가 $0.99$ 근처에서 $\pm0.005$만큼 무작위로 변한다. 간격은 어느 범위에 있는가?
\end{exercise}

\begin{exercise}[\lvlA]
\label{ex:21:2}
\emph{테스트 벤치의 에너지.} (a)~추정~\eqref{eq:dishpower}은 뇌 전체($8.6\cdot10^{10}$개 뉴런)에 대해 얼마의 전력을 주는가? (b)~테스트 벤치가 전극 $1024$개를 $20$~kHz로, 표본당 $10$~비트로 기록한다. 비트당 $1$~nJ일 때 이 데이터 흐름의 전송은 배양 자체의 전력보다 몇 배 비싼가?
\end{exercise}

\begin{exercise}[\lvlB]
\label{ex:21:3}
\emph{설계: 일제 발화 억제.} 여러 전극을 통한 드문 자극 때문에 모든 시냅스가 빈도 $r_b$로 신경전달물질을 방출한다. 저장량은 $\tau_{\text{rec}}\,\dd x/\dd t=1-x-Ur_b\tau_{\text{rec}}x$, $U=0.5$, $\tau_{\text{rec}}=0.8$~s. 저장량이 $x^*=0.9$; $0.99$에 이르지 못하게 하는 가장 작은 $r_b$는 얼마이며, 그때 시냅스는 얼마나 약해지는가?
\end{exercise}

\begin{exercise}[\lvlB]
\label{ex:21:4}
\emph{저수지의 기억은 뉴런 수를 넘지 않는다.} 출력이 $N$개인 저수지 $\mathbf r(t)$가 평균 0, 분산 1인 백색 입력 $u(t)$를 받는다. $m_k$는 판독 $\mathbf w_k\cdot\mathbf r(t)$와 $u(t-k)$의 상관계수 제곱의 최댓값이다. 기억 용량 $\mathrm{MC}=\sum_{k\ge0}m_k\le N$임을 증명하라.
\end{exercise}

\begin{exercise}[\lvlB]
\label{ex:21:5}
\emph{모의실험: 실리콘 저수지.} 저수지 $\mathbf r(t+1)=\tanh(W\mathbf r(t)+\mathbf v\,u(t)+\mathbf c)$, $N=30$, $W$는 스펙트럼 반지름 $0.9$인 무작위 행렬, $v_i\sim U(-0.5,0.5)$, $c_i\sim U(-0.2,0.2)$, $u\sim U(-1,1)$, 판독은 릿지 회귀이다. $\tanh$가 있을 때와 없을 때 연습문제~\ref{ex:21:4}의 기억 용량을 구하라. 어느 저수지가 더 많이 기억하는가?
\end{exercise}

\begin{exercise}[\lvlB]
\label{ex:21:6}
\emph{살아 있는 저수지의 판독.} 오가노이드가 채널 $1024$개와 두 클래스의 학습 기록 $P=240$개를 준다. (a)~연습문제~\ref{ex:19:3}의 공식으로 무작위 라벨이 분리될 확률이 $1\%$ 이하가 되려면 특징을 몇 개($n$)까지 남길 수 있는가? (b)~검증 기록 $100$개에서 정확도 $78\%$는 찍기보다 몇 시그마 높은가?
\end{exercise}

\begin{exercise}[\lvlC]
\label{ex:21:7}
\emph{연구: 브로드캐스트 채널 없는 교사.} 전극 $8$--$1024$개로 테스트 벤치에서 \ref{sec:dofa}장의 3요소 규칙을 구현하는 자극 프로토콜과, 가중치 학습을 네트워크 상태의 이동과 구별하는 판독 고정 대조 실험을 제안하라. 어떤 결과가 학습 가설을 반증하겠는가?
\end{exercise}
